\section{Introduction}
\label{sec:introduction}

\subsection{The Puzzle and Promise of In-Context Learning}
\label{subsec:background}

In-context learning (ICL) stands as one of the most remarkable and somewhat mysterious capabilities of modern large language models (LLMs) \cite{brown2020language}. Unlike traditional fine-tuning, ICL allows a model to adapt to new tasks solely through inference-time conditioning on a few input-output examples provided in the prompt, without any parameter updates. This meta-learning-like behavior emerges unexpectedly during pre-training on next-token prediction and has become a cornerstone of practical LLM deployment.

The theoretical community has made significant strides in demystifying this phenomenon, particularly for \textit{linear} transformer architectures. Seminal works by \cite{von2023linear}, \cite{akyurek2022learning}, and others have elegantly shown that linear attention heads can implement gradient descent on a least-squares objective, making ICL equivalent to solving ridge regression. This line of work paints a compelling picture: ICL in linear models is essentially implicit Bayesian inference or gradient-based optimization performed in the forward pass.

However, reality is nonlinear. Every transformer block that powers today's LLMs contains two fundamentally nonlinear components: the \textbf{softmax attention} mechanism and the \textbf{multi-layer perceptron (MLP)} with nonlinear activations like ReLU or GeLU. While the linear theory is beautiful and instructive, it necessarily ignores these core architectural elements. We are thus left with a critical gap: if ICL in linear models is ridge regression, what is ICL in real, nonlinear transformers? What kind of learning algorithm does it implicitly execute, and what statistical guarantees can we derive?

Understanding the implicit regularization induced by these nonlinearities is not merely an academic exercise. It is the key bridge connecting the elegant but simplified linear theories to the messy, powerful reality of practical models. This understanding could illuminate why certain prompts work better than others, guide architecture design, and ultimately help us build more reliable and efficient in-context learners.

\subsection{The Limits of Linearity and Our Position}
\label{subsec:positioning}

Existing theoretical work provides essential pieces, but not the complete puzzle for nonlinear ICL.

\textbf{Linear ICL Theory.} The equivalence between linear transformers and gradient descent/ridge regression is a foundational result \cite{von2023linear, akyurek2022learning, zhang2023trained}. These analyses, however, rely crucially on the assumption of linear attention or linearized approximations. They cannot capture the expressive power or inductive biases introduced by nonlinear activations. As noted in \cite{huang2023understanding}, the linear theory fundamentally cannot explain how models perform ICL on nonlinear tasks.

\textbf{Classical Implicit Regularization.} There is a rich literature on how gradient-based optimization imparts implicit biases, such as convergence to max-margin solutions in linear models \cite{soudry2018implicit} or low-rank solutions in matrix factorization \cite{gunasekar2018implicit}. The Neural Tangent Kernel (NTK) theory \cite{jacot2018neural, arora2019exact} further shows that infinite-width networks trained with gradient descent behave like kernel regression. However, these analyses typically focus on a \textit{fixed} dataset and a \textit{single} task. ICL, in stark contrast, is a meta-learning scenario: the model is trained on a distribution of tasks, and at test time, the "dataset" (the in-context examples) is dynamic and task-specific. The object of regularization is not a single function but the learning algorithm itself.

\textbf{Attention as Kernel Methods.} The connection between softmax attention and kernel smoothing has been explored \cite{tsai2019transformer, katharopoulos2020transformers}, and recent work by \cite{chen2024kernels} has framed ICL as kernel gradient descent in a feature space defined by attention. Yet, these works often treat the attention kernel in isolation or within linearized settings. They do not provide a unified analysis of the \textit{full transformer block}—where attention and MLP interact—and its consequent implicit regularization.

\textbf{Emerging Explorations of Nonlinearity.} Recent studies have begun to probe the necessity of nonlinearity. For instance, \cite{sohn2024understanding} empirically demonstrate that the MLP is crucial for nonlinear ICL, while linearized transformers fail. \cite{liu2023theoretical} provide initial steps by analyzing shallow ReLU networks in an ICL setting. However, a \textit{systematic theoretical framework} that characterizes the implicit regularization of standard transformer architectures (softmax attention + nonlinear MLP) and derives finite-sample generalization bounds remains absent.

This paper aims to fill this exact gap. We move \textbf{beyond the linear paradigm} to establish the first comprehensive theory of implicit regularization in nonlinear in-context learning.

\subsection{Our Contributions}
\label{subsec:contributions}

We present a unified theoretical framework that characterizes how a single-layer transformer with softmax attention and a nonlinear MLP performs ICL. Our main contributions are fourfold:

\begin{enumerate}
	\item \textbf{A Kernel Representation Theorem for Nonlinear ICL.} We prove that gradient flow training of a transformer induces a predictor that solves a \textit{kernel ridge regression} (KRR) problem. Crucially, the kernel is not arbitrary; it is a structured combination $\alpha K_{\text{att}} + \beta K_{\text{mlp}}$, where $K_{\text{att}}$ is an exponential kernel induced by the attention mechanism and $K_{\text{mlp}}$ is the ReLU neural tangent kernel (Theorem 1). This provides a precise functional characterization of the nonlinear ICL algorithm.
	
	text
	\item \textbf{Tight Generalization Bounds.} We derive non-asymptotic, high-probability upper bounds on the excess risk of the in-context learner (Theorem 2). These bounds explicitly depend on interpretable quantities:
	\begin{itemize}
		\item The \textit{effective dimension} $d_{\text{eff}}$ of the combined kernel, capturing task complexity.
		\item The Lipschitz constant of the activation function (e.g., 1 for ReLU).
		\item The number of in-context examples $n$.
	\end{itemize}
	The bound decays as $O(\sqrt{d_{\text{eff}}/n})$, providing a rigorous sample complexity guarantee for nonlinear ICL.
	
	\item \textbf{Discovery of a Regularization Phase Transition.} We identify a surprising \textbf{phase transition} in the implicit regularization mechanism (Theorem 3). There exists a critical context length $n_c \propto d_{\text{eff}}$. When $n < n_c$ (under-specified regime), the implicit bias resembles an $\ell_2$-norm penalty, promoting smooth solutions. When $n > n_c$ (over-specified regime), the bias morphs into an $\ell_1$-type penalty, encouraging \textit{sparse} representations in the kernel basis. This phenomenon explains how transformers can adapt their inductive bias based on available context.
	
	\item \textbf{Empirical Validation and Practical Insights.} We validate all theoretical predictions on synthetic tasks (polynomial regression, RBF functions). The empirical risk curves closely match our derived bounds, and we clearly observe the predicted phase transition. Our theory yields actionable insights:
	\begin{itemize}
		\item \textbf{Prompt Design:} For simple tasks (low $d_{\text{eff}}$), short prompts suffice. For complex tasks, prompts must exceed the critical length $n_c$ to activate the beneficial sparse regime.
		\item \textbf{Architecture Guidance:} The width of the MLP controls the richness of $K_{\text{mlp}}$, while the number of attention heads diversifies $K_{\text{att}}$.
		\item \textbf{Meta-Training:} The task distribution during pre-training shapes the effective kernel and thus the sample efficiency on downstream tasks.
	\end{itemize}
\end{enumerate}

In essence, this paper constructs a bridge from the well-understood land of linear ICL to the complex continent of nonlinear transformers. We show that the path across this bridge is illuminated by kernel methods and spectral analysis, revealing a landscape richer and more dynamic than previously imagined—one where the learning algorithm itself evolves with the length of the prompt.

\subsection{Paper Organization}
The rest of the paper is structured as follows. Section 2 reviews related work. Section 3 formalizes our problem setup and model architecture. Our main theoretical results—Theorems 1, 2, and 3—are presented in Section 4, with proof sketches. Detailed technical analyses and full proofs are deferred to Sections 5 and the Appendix. Section 6 presents empirical validations. We discuss broader implications, limitations, and future directions in Section 7, and conclude in Section 8.